\subsection{Convex cones in topological vector spaces}

PROPOSITION 14. --- In a topological vector space E, the closure of a convex set (resp. of a convex cone) is a convex set (resp. a convex cone with the same vertex).

For, let A be a convex set; the mapping \((x, y) \mapsto \lambda x + (1 - \lambda)y\) , where \(0 < \lambda < 1\) , is continuous in E × E and maps A × A in A; thus (GT, I, §2.1, th. 1) it maps \(\overline{A} \times \overline{A}\) in \(\overline{A}\) , which shows that \(\overline{A}\) is convex. Similarly, if C is a convex cone with vertex 0 then \(\overline{C} + \overline{C} \subset \overline{C}\) and \(\lambda\overline{C} \subset \overline{C}\) for all \(\lambda > 0\) .

DEFINITION 4. --- For any set A of a topological vector space E, the intersection of all the closed convex sets containing A is called the convex closed envelope of A; it is the smallest convex closed set containing A.

From prop. 14, the convex closed envelope of A is the closure of the convex envelope of A; it is clearly the same as the convex closed envelope of \(\overline{A}\) .

Similarly we call the smallest symmetric, convex, closed set that contains A, the symmetric convex closed envelope (or the balanced convex closed envelope) of A; it is the closure of the symmetric convex envelope of A (II, p.~10); it is also the symmetric convex closed envelope of \(\overline{A}\) .

PROPOSITION 15. --- Let \(A_{i}\) ( \(1 \leqslant i \leqslant n\) ) be a finite number of compact convex sets in a Hausdorff topological vector space E. Then the convex envelope of the union of the \(A_{i}\) is compact (and is, therefore, the same as the convex closed envelope of this union).

Let B be the compact set in \(R^{n}\) defined by the points \((\lambda_{1},\lambda_{2},\ldots,\lambda_{n})\) where \(\lambda_{i}\geqslant0(1\leqslant i\leqslant n)\) , and \(\sum_{i=1}^{n}\lambda_{i}=1\) . Define a continuous mapping of \(B\times\prod_{i=1}^{n}A_{i}\subset R^{n}\times E^{n}\) in E by the formula \((\lambda_{1},\lambda_{2},\ldots,\lambda_{n},x_{1},x_{2},\ldots,x_{n})\mapsto\sum_{i=1}^{n}\lambda_{i}x_{i}\) . The convex envelope C of \(\bigcup_{i=1}^{n}A_{i}\) is the image of \(B\times\prod_{i=1}^{n}A_{i}\) under this mapping; as \(B\times\prod_{i=1}^{n}A_{i}\) is compact and E is Hausdorff, it follows that C is compact.

COROLLARY 1. --- In a Hausdorff topological vector space the convex envelope of a finite set is compact.

COROLLARY 2. --- In a topological vector space E, the convex envelope of a finite set is precompact.

In fact, let \(j\) be the canonical mapping of \(\mathbf{E}\) in its Hausdorff completion \(\hat{\mathbf{E}}\) ; if \(\mathbf{C}\) is the convex envelope of \(\mathbf{A}\) , then \(j(\mathbf{C})\) is the convex envelope of the finite set \(j(\mathbf{A})\) in \(\hat{\mathbf{E}}\) , hence \(j(\mathbf{C})\) is compact (cor. 1) and therefore \(\mathbf{C}\) is precompact (GT, II, § 4.2).

PROPOSITION 16. --- Let A be a convex subset, with at least one interior point \(x_{0}\) , of a topological vector space E. For any point \(x \in \overline{A}\) , every point of the open segment with end points \(x_{0}\) , x lies in the interior of A.

For any point y of this segment, let f be the homothety of centre y and ratio \(\lambda < 0\) , which transforms \(x_{0}\) into x. If V is an open neighbourhood of \(x_{0}\) contained in A, then \(f(\mathrm{V})\) is a neighbourhood of x and therefore contains a point \(f(z) \in \mathrm{A}\) ; now

\[
f (z) - y = \lambda (z - y) = \lambda (z - f (z)) + \lambda (f (z) - y),
\]

hence \(y - f(z) = \frac{\lambda}{\lambda - 1} (z - f(z))\) , so that \(y\) is transformed into \(z\) by the homothety \(g\) , of centre \(f(z)\) and ratio \(\mu = \lambda / (\lambda - 1)\) ; since \(0 < \mu < 1\) , \(g\) transforms \(V\) into a neighbourhood of 0 contained in \(A\) . The proposition is proved.

COROLLARY 1. --- The interior \(\mathring{\mathrm{A}}\) of a convex set \(\mathrm{A}\) , is itself a convex set; if \(\mathring{\mathrm{A}}\) is not empty, then it coincides with the interior of \(\overline{\mathrm{A}}\) , and \(\overline{\mathrm{A}}\) is a convex set that coincides with the closure of \(\mathring{\mathrm{A}}\) .

It follows from prop. 16, that if \(\mathring{\mathrm{A}}\) is not empty, then it is a convex set and every point of \(\overline{\mathrm{A}}\) is a cluster point of \(\mathring{\mathrm{A}}\) . Next we show that every interior point of \(\overline{\mathrm{A}}\) belongs to \(\mathring{\mathrm{A}}\) . Let \(x\) be an interior point of \(\overline{\mathrm{A}}\) and suppose, for definiteness that \(x = 0\) . Let \(\mathrm{V}\) be a symmetric neighbourhood of 0 that is contained in \(\overline{\mathrm{A}}\) and let \(y \in \mathring{\mathrm{A}} \cap \mathrm{V}\) ; now \(-y \in \overline{\mathrm{A}}\) , and therefore, by prop. 16, we see that \(0 \in \mathring{\mathrm{A}}\) , if \(y \neq 0\) ; this is obviously true if \(y = 0\) .

COROLLARY 2. --- The interior \(\mathring{\mathrm{C}}\) of a convex cone \(\mathrm{C}\) , is itself a convex cone; if \(\mathring{\mathrm{C}}\) is not empty then it coincides with the interior of \(\overline{\mathrm{C}}\) , and \(\overline{\mathrm{C}}\) is a pointed convex cone that coincides with the closure of \(\mathring{\mathrm{C}}\) .

Since homotheties of ratio \textgreater{} 0 and centre 0 transform C into itself, they do the same for \(\overset{\circ}{C}\) , thus \(\overset{\circ}{C}\) is a cone; the remainder of the corollary follows from cor. 1 and the obvious remark that if C is not empty then \(\overline{C}\) contains the vertex of C.

Let H be a closed hyperplane in the topological vector space E over R; it has an equation of the form \(f(x) = \alpha\) , where f is a continuous linear form that is not identically zero in E (I, p.~13, th. 1). The closed half spaces defined respectively by \(f(x) \leqslant \alpha\) and \(f(x) \geqslant \alpha\) are therefore closed convex sets; their complements defined respectively by \(f(x) > \alpha\) and \(f(x) < \alpha\) , are open convex sets. We say that these half-spaces are the closed (resp. open) half spaces determined by H.

PROPOSITION 17. --- In a topological vector space E, let A be a set with at least one interior point, and such that all its points lie on the same side of an hyperplane H. Then H is closed, the interior points of A lie strictly on the same side of H, and the cluster points of A lie on the same side of H. In particular open (resp. closed) half spaces are determined by closed hyperplanes.

In fact suppose that H contains the origin and that \(f(x) = 0\) is an equation of H; suppose, for definiteness, that \(f(x) \geqslant 0\) for all \(x \in A\) . The half space formed by the points y such that \(f(y) > -1\) contains at least one interior point, and, by translation, the same is true of the half space of points y such that \(f(y) > 0\) ; this shows that H is closed (I, p.~11, corollary). Then we know that f is a strict morphism of E on R (I, p.~13, corollary), therefore \(f(\mathring{\mathrm{A}})\) is an open set in R. This set cannot contain 0 or it would contain numbers \textless{} 0 contrary to hypothesis; it is thus contained in the open interval \(\bigg]0, +\infty[\) . On the other hand, the half space of those y for which \(f(y) \geqslant 0\) is closed and contains A, therefore it contains \(\overline{A}\) .

COROLLARY. --- Let P be a pointed convex cone, with at least one interior point, of the topological vector space E. Then each linear form f that is not identically zero on E, and is positive for the preorder structure defined by P(II, p.~13), is necessarily continuous. Further, if x is interior to P then \(f(x) > 0\) and if x is a cluster point of P then \(f(x) \geqslant 0\) .

Apply prop. 17 to the case \(\mathbf{A} = \mathbf{P}\) where \(\mathbf{H}\) is the hyperplane with the equation \(f(x) = 0\) .

Remark. --- In a topological vector space E, every convex set C is connected. In fact, if \(a \in C\) , then C is a union of segments with end point a and closed at a; these are connected and the result follows from GT, I, § 11.1, prop. 2.

